% ============================================================================
% RMG-DFT TDDFT refactor: old-to-new source mapping (May 2026 upstream sync)
% ============================================================================
% Compile: pdflatex tddft_refactor_mapping
% ============================================================================
\documentclass[11pt,letterpaper]{article}

\usepackage{amsmath,amssymb}
\usepackage{hyperref}
\usepackage{booktabs}
\usepackage{longtable}
\usepackage{listings}
\usepackage{xcolor}
\usepackage{geometry}
\usepackage{enumitem}
\usepackage{tcolorbox}
\usepackage{bm}

\geometry{margin=1in}

\hypersetup{
  colorlinks=true,
  linkcolor=blue!70!black,
  citecolor=green!50!black,
  urlcolor=blue!70!black
}

% ---------- C++ listing style (matches rmg_tddft_theory.tex) ----------
\lstdefinestyle{cppstyle}{
  language=C++,
  basicstyle=\small\ttfamily,
  keywordstyle=\color{blue!80!black}\bfseries,
  commentstyle=\color{green!50!black}\itshape,
  stringstyle=\color{orange!80!black},
  frame=single,
  framerule=0.4pt,
  rulecolor=\color{gray!50},
  backgroundcolor=\color{gray!5},
  breaklines=true,
  showstringspaces=false,
  tabsize=4
}

% ---------- Math macros (shared with rmg_tddft_theory.tex) ----------
\newcommand{\ii}{\mathrm{i}}

% ---------- Boxes for gaps and new items ----------
\tcbuselibrary{skins,breakable}
\newtcolorbox{gapbox}[2][]{
  title={#2},
  colback=red!5!white,
  colframe=red!50!black,
  fonttitle=\bfseries\small,
  breakable,
  #1
}
\newtcolorbox{infobox}[2][]{
  title={#2},
  colback=blue!5!white,
  colframe=blue!50!black,
  fonttitle=\bfseries\small,
  breakable,
  #1
}

% ============================================================================
\begin{document}

\title{\Large\bfseries TDDFT Refactor: Source-Code Mapping\\[4pt]
\normalsize\textit{Pre-refactor \texttt{RmgTddft.cpp} $\to$
Post-refactor \texttt{rmg::tddft} class}\\[2pt]
\normalsize\textit{Companion to \texttt{rmg\_tddft\_theory.tex}}}
\author{Jacek Jakowski \and \textit{(working notes, May 2026)}}
\date{\today}
\maketitle

\begin{abstract}
In May 2026 the RMG-DFT upstream repository underwent a structural refactor
that affected the TDDFT module.  The pre-refactor monolithic free function
\texttt{RmgTddft<O,M>(...)} was deleted; its functionality was redistributed
into a new \texttt{rmg::tddft<OrbitalType,MatrixType>} class split across
three source files.  The annotated copy of the old file in the frozen
\texttt{explore} branch serves as a \emph{Rosetta Stone} for understanding
what the new code does and where each piece of the old algorithm now lives.

This note documents the structural change, the section-by-section mapping
between old and new code, the post-refactor status of the eight known
theory/code gaps (\texttt{GAP-1}$\dots$\texttt{GAP-8}) catalogued in the
companion theory document, and three substantive additions that the
refactor introduced.  It serves as an entry point for further work on
non-collinear spin RT-TDDFT and Ehrenfest dynamics.
\end{abstract}

\tableofcontents
\bigskip

% ============================================================================
\section{Context: why this document exists}
\label{sec:context}
% ============================================================================

The \texttt{explore} branch of the personal fork
\texttt{jjakowski/rmgdft-sandbox} was used for several months as a sandbox
for annotating the RT-TDDFT C++ source with detailed physics commentary.
Three files were annotated:
\texttt{TDDFT/RMG\_TDDFT/RmgTddft.cpp}, \texttt{CurrentNlpp.cpp}, and
\texttt{VecPmatrix.cpp}.  The manuscript and theory notes also live on
\texttt{explore}.

When the upstream sync was performed in May 2026, three issues appeared:
\begin{itemize}[leftmargin=*]
  \item \texttt{RmgTddft.cpp} was \emph{deleted} from upstream and replaced
        by a different file structure --- annotation could not be carried
        forward by simple copy.
  \item \texttt{CurrentNlpp.cpp} and \texttt{VecPmatrix.cpp} were unchanged
        upstream --- annotation \emph{could} be copied verbatim.
  \item \texttt{explore} also held months of manuscript work that must be
        preserved.
\end{itemize}
The resolution was to (i) tag the pre-sync state of \texttt{develop} as
\texttt{pre-TDDFT-annotation}, (ii) leave \texttt{explore} frozen as a
permanent Rosetta Stone, (iii) sync \texttt{develop} with upstream, and
(iv) create a clean \texttt{explore-new} branch off the post-sync
\texttt{develop}.  Annotations from the unchanged files were migrated by
direct copy (Task~A).  Annotations for \texttt{RmgTddft.cpp} were migrated
by re-mapping each annotated section to its new location in the refactored
files (Task~B).  This document captures the structural understanding
developed during Task~B.

% ============================================================================
\section{Structural change}
\label{sec:structure}
% ============================================================================

\paragraph{Before.}  A single templated free function did everything
inline:
\begin{lstlisting}[style=cppstyle]
template <typename OrbitalType, typename MatrixType>
void RmgTddft( spinobj<double>& vxc,
               fgobj<double>&   vh,
               fgobj<double>&   vnuc,
               spinobj<double>& rho,
               fgobj<double>&   rhocore,
               fgobj<double>&   rhoc,
               Kpoint<OrbitalType>** Kptr );
\end{lstlisting}
1245 lines in one function: SETUP (memory, file open, \texttt{ReadData})
$\to$ velocity-gauge init / kick $\to$ ground-state $J_0$ baseline $\to$
outer time loop (STEP 1 / 2 / 3) $\to$ cleanup (file close, write
checkpoint).

\paragraph{After.}  A C++ class
\texttt{rmg::tddft<OrbitalType,MatrixType>} split across three files:
\begin{center}
\begin{tabular}{ll}
\toprule
File & Members \\
\midrule
\texttt{rmg\_tddft.cpp} & constructor \texttt{tddft(...)} (setup + init + kick) \\
                       & method \texttt{tddft\_md()} (outer time loop) \\
                       & destructor \texttt{$\sim$tddft()} (write + close + free) \\
                       & helpers \texttt{tstconv()}, \texttt{gather\_rho\_matrix()} \\
\texttt{rmg\_tddft\_energy.cpp} & methods \texttt{tddft\_energy\_init()}, \texttt{tddft\_energy()} \\
\texttt{rmg\_rotate\_sint.cpp}  & free function \texttt{rmg::rotate\_sint<O>()} (Ehrenfest helper) \\
\bottomrule
\end{tabular}
\end{center}

The split reflects a real architectural change, not just file
reorganisation: in the new design \texttt{tddft\_md()} can be invoked
\emph{multiple times} from an outer Ehrenfest MD loop, each invocation
propagating the electrons for \texttt{ct.tddft\_steps} steps before
returning forces on the nuclei.  The constructor sets up the persistent
state once; the destructor flushes it once.  In the old design every TDDFT
call was a single self-contained run from setup to cleanup, with no notion
of repeated invocations.

% ============================================================================
\section{Section-by-section mapping}
\label{sec:mapping}
% ============================================================================

Line numbers in the rightmost column refer to
\texttt{TDDFT/RMG\_TDDFT/rmg\_tddft.cpp} after the Task~B annotation commit
(\texttt{a998da7c2}).  Line numbers in the leftmost column refer to the
frozen \texttt{explore} version of \texttt{RmgTddft.cpp}.

\begin{longtable}{p{0.32\linewidth} p{0.20\linewidth} p{0.40\linewidth}}
\toprule
\textbf{Old section (RmgTddft.cpp)} & \textbf{New member} & \textbf{New location (rmg\_tddft.cpp)} \\
\midrule
\endhead
File-level header + GAP-1..8 audit (l.~23--152)
  & --- (file header)
  & Top of file, after GPL header (l.~25--165) \\
\midrule
Velocity-gauge initialization (l.~606--657)
  & constructor
  & Inside \texttt{if(ct.tddft\_mode == VECTOR\_POT)} block, l.~561 \\
Inline \texttt{STEP A / B / C} labels (l.~663, 668, 677)
  & constructor
  & l.~615, 620, 629 \\
Ground-state $J_0$ baseline (l.~696--707)
  & constructor
  & l.~648 (before zdotc loop populating \texttt{current0[]}) \\
\midrule
Outer time loop comment (l.~794--803)
  & \texttt{tddft\_md()}
  & l.~830 (before \texttt{for(int tddft\_steps...)}) \\
STEP 1 --- PREDICTOR (l.~815--820)
  & \texttt{tddft\_md()}
  & l.~851 (before \texttt{extrapolate\_Hmatrix}) \\
STEP 2 --- CORRECTOR SCF (l.~849--854)
  & \texttt{tddft\_md()}
  & l.~886 (before SCF \texttt{while}) \\
STEP 2a --- Magnus operator (l.~885--889)
  & \texttt{tddft\_md()}
  & l.~921 (before \texttt{magnus(...)}) \\
STEP 2b --- BCH propagator (l.~892--898)
  & \texttt{tddft\_md()}
  & l.~929 (before \texttt{eldyn\_ort(...)}) \\
STEP 2c --- Density update (l.~916--920)
  & \texttt{tddft\_md()}
  & l.~955 (before Rho \texttt{RmgTimer}) \\
STEP 2d --- Update potentials (l.~971--973)
  & \texttt{tddft\_md()}
  & l.~1009 (before \texttt{compute\_vxc(...)}) \\
STEP 2e --- Update Hamiltonian (l.~991--999)
  & \texttt{tddft\_md()}
  & l.~1029 (before \texttt{vtot} accumulation) \\
Convergence check (l.~1038--1042)
  & \texttt{tddft\_md()}
  & l.~1078 (after \texttt{daxpy}, before \texttt{tstconv}) \\
STEP 3 --- Advance (l.~1096--1101)
  & \texttt{tddft\_md()}
  & l.~1144 (before advance \texttt{for} loop) \\
TIME-DEPENDENT CURRENT $J(t)$ (l.~1113--1130)
  & \texttt{tddft\_md()}
  & l.~1160 (inside \texttt{if(VECTOR\_POT)} branch) \\
\bottomrule
\end{longtable}

\noindent\textit{All twelve annotation blocks were ported verbatim from the
Rosetta Stone}; only their position changed to follow the constructor /
\texttt{tddft\_md()} split.

% ============================================================================
\section{GAP audit (post-refactor)}
\label{sec:gaps}
% ============================================================================

Each of the eight gaps documented in
\href{rmg_tddft_theory.tex}{\texttt{rmg\_tddft\_theory.tex}} was re-checked
against the refactored code.  \emph{None was fixed by the refactor.}
The evidence for each is summarised below.

\begin{gapbox}{GAP-1: Velocity-gauge mode is delta-kick only}
The time-varying field update
\texttt{daxpy(A\_0*cos($\omega t$), Pxmatrix, Hmatrix\_0)}
is still present but \emph{commented out} (\texttt{rmg\_tddft.cpp} l.~838--846
inside \texttt{tddft\_md()}).  \texttt{ct.tddft\_frequency} is read from input
but has no effect on the propagation.
\end{gapbox}

\begin{gapbox}{GAP-2: Diamagnetic current is missing}
\texttt{current[]} is still computed as $J_\alpha = \mathrm{Re}\,
\mathrm{Tr}(P\, P^\alpha)$ only --- see l.~1175--1186 inside
\texttt{tddft\_md()}.  The $-n(r,t)\,A(t)$ contribution is absent.  For the
delta-kick (where $A \to 0$ after $t=0$) this is fine; if GAP-1 is fixed,
GAP-2 must also be.
\end{gapbox}

\begin{gapbox}{GAP-3: \texorpdfstring{$A^2$}{A^2} term is dropped}
The minimal-coupling Hamiltonian $(p + A)^2/2 = p^2/2 + A\cdot p + A^2/2$
is truncated to the linear $A\cdot p$ term.  No amplitude check enforces the
linear-response regime in which this is valid.
\end{gapbox}

\begin{gapbox}{GAP-4: Non-collinear spin in VECTOR\_POT path is incomplete}
The \texttt{vtot} update at STEP 2e uses only the scalar part of $V_{xc}$.
Same code structure as before; both \texttt{// noncoll need change} comments
are still present at l.~1039 and l.~1048 of \texttt{rmg\_tddft.cpp}.
This is directly relevant to the planned non-collinear spin RT-TDDFT work.
\end{gapbox}

\begin{gapbox}{GAP-5: \texttt{Hmatrix\_m1\_cpu} triple-purposed}
Same buffer reuse pattern as before: before SCF it is $H(t-dt)$; during
STEP 2a it is overwritten with the Magnus operator $\Omega$; during STEP 2e
it is overwritten with $\langle \delta V \rangle$.  $H(t-dt)$ is restored
only at STEP 3 via \texttt{memcpy(Hmatrix\_m1\_cpu, Hmatrix\_0\_cpu, ...)}.
\end{gapbox}

\begin{gapbox}{GAP-6: \texttt{Hmatrix\_cpu} is never explicitly reset to $H(t)$}
Same convention as before: \texttt{Hmatrix\_cpu} carries the converged $H$
from the previous step and the SCF loop adds \emph{incremental} potential
corrections $\delta V$ on top of it.  Any future modification that breaks
the ``$vtot$ is incremental'' invariant will silently corrupt $H$.
\end{gapbox}

\begin{gapbox}{GAP-7: No guard on mode vs.~boundary conditions}
No check that EFIELD is used only with finite molecules
(\texttt{ct.is\_gamma == true}) or that VECTOR\_POT is used only with
proper k-point sampling.  Misuse compiles and runs but produces
unphysical results.
\end{gapbox}

\begin{gapbox}{GAP-8: \texttt{ct.efield\_tddft\_crds} semantics differ between modes}
The same input array is interpreted as $E_0$ (length gauge, EFIELD) or
$A_0$ (velocity gauge, VECTOR\_POT).  The relationship $E_0 = A_0 \omega$
means the two modes require different input values for the same physical
field strength.  Not enforced in input parsing.
\end{gapbox}

\noindent\textbf{Implication.}  The refactor was a \emph{structural}
reorganisation, not a physics rewrite.  Every theoretical limitation
identified before the refactor remains an open item in
\texttt{explore-new}.

% ============================================================================
\section{New items introduced by the refactor}
\label{sec:newitems}
% ============================================================================

Three blocks of code have no analogue in the old \texttt{RmgTddft.cpp}.
They are documented inline in the post-refactor source; the discussion
below is for forward reference.

% ----------------------------------------------------------------------------
\subsection{Per-call bootstrap (\texttt{first\_step})}
\label{sec:bootstrap}

Located near the top of \texttt{tddft\_md()} (l.~776--800):
\begin{lstlisting}[style=cppstyle]
static int first_step;
for(int kpt = 0; kpt < ct.num_kpts_pe; kpt++)
{
    BetaProjector->project(...);    // ALWAYS: refresh sint = <beta|phi>
    if(first_step) {
        HSmatrix(...);              // rebuild H_ij = <phi_i|H_KS|phi_j>
        DistributeMatrix(...);      // scatter into BLACS storage
        // re-seed Hmatrix_m1/0/1 with fresh H
    }
}
first_step++;
\end{lstlisting}

The condition \texttt{if(first\_step)} evaluates to \texttt{false} on the
very first call (since the \texttt{static int} is zero-initialized) and
\texttt{true} on every subsequent call.  This skips the $H$ rebuild on the
first invocation --- because the constructor already filled $H$ with the
diagonal of ground-state eigenvalues, which is valid as long as ions have
not moved.  From the second call onward (i.e., after the outer MD loop
has moved ions), the ground-state diagonal $H$ is no longer valid and
must be rebuilt from scratch via \texttt{HSmatrix} in the current geometry.

\begin{infobox}{Implication for non-collinear spin RT-TDDFT}
For non-collinear spinor wavefunctions the bootstrap path matters because
\texttt{HSmatrix} must produce the correct $2\times 2$ spin-space matrix
elements.  Verify that the \texttt{HSmatrix} implementation in
\texttt{TDDFT/RMG\_TDDFT/HSmatrix.cpp} respects \texttt{ct.noncoll\_factor}
before enabling noncoll with multi-step Ehrenfest dynamics.
\end{infobox}

% ----------------------------------------------------------------------------
\subsection{Running time-averaged density (\texttt{trho})}
\label{sec:trho}

After each propagation step (\texttt{rmg\_tddft.cpp} l.~1104--1111):
\[
  \texttt{trho}(\mathbf{r}) \mathrel{+}= \rho(\mathbf{r}, t_n)
\]
After the outer time loop completes, \texttt{trho} is rescaled by
$1/N_\text{steps}$:
\[
  \overline{\rho}(\mathbf{r})
  = \frac{1}{N_\text{steps}}\sum_{n=0}^{N_\text{steps}-1} \rho(\mathbf{r},t_n).
\]
The time-averaged density $\overline\rho$ is the input to the
\texttt{Force(...)} call in the Ehrenfest block (Section~\ref{sec:cve}),
rather than the final-snapshot density $\rho(t_\text{final})$.

\paragraph{Why averaging?}  The instantaneous TD density carries fast
electronic oscillations (interband and plasmon-like) that, if used
directly, would produce noisy nuclear forces and unstable trajectories.
Averaging over the propagation window damps these oscillations and
provides smoother forces.  This is conceptually similar to the
``time-averaged forces'' used in adiabatic-following schemes in
quantum-chemistry molecular dynamics codes; physically it is a low-pass
filter on $F[\rho(t)]$.

% ----------------------------------------------------------------------------
\subsection{Ehrenfest force update (\texttt{TDDFT\_CVE} branch)}
\label{sec:cve}

The post-refactor \texttt{tddft\_md()} ends with an Ehrenfest force-update
block (l.~1259--1305 in \texttt{rmg\_tddft.cpp}).  It runs when
\texttt{ct.forceflag == TDDFT\_CVE}, i.e., when the user has requested
Combined-Velocity-Verlet/Ehrenfest dynamics.  The algorithm:

\begin{enumerate}[leftmargin=*]
\item \textbf{Rotate force pointers.}  \texttt{Atoms[ion].RotateForces()}
      shifts old per-ion force snapshots back by one slot so the
      Beeman/Velocity-Verlet integrator has access to $F(t-dt)$,
      $F(t-2dt)$ after the new force is written.
\item \textbf{Gather global density matrix.}
      \texttt{gather\_rho\_matrix(global, Pn0)} collects the
      BLACS-distributed $P(t_\text{final})$ into a global
      $N_\text{states} \times N_\text{states}$ matrix on each rank.
\item \textbf{Rotate \texttt{sint} into TD basis.}
      \texttt{save\_sint()} preserves the ground-state
      $\texttt{sint} = \langle \beta_{lm} | \phi_j \rangle$;
      occupations are temporarily set to 1.0 (slot~[3] holds the
      ground-state values for restore);
      \texttt{rmg::rotate\_sint(...)} multiplies the ground-state
      \texttt{sint} by the density matrix to produce the TD-basis
      \texttt{sint}.  Occupations $f_j$ are set to unity because the
      per-orbital weights are already encoded in $P_{ij}$.
\item \textbf{Compute forces from time-averaged density.}
      \texttt{Force(trho.up, trho.dw, ..., Kptr)} computes the nuclear
      forces using $\overline{\rho}$, not the final snapshot.
\item \textbf{Restore state.}  \texttt{restore\_sint()}, restore
      occupations, then \texttt{get\_ddd(...)} refreshes the $D_{ij}$
      dual-projector coefficients for the next call.
\end{enumerate}

\begin{infobox}{Open: nonlocal-PP contribution to the Ehrenfest force}
The current \texttt{Force(...)} call assumes the local potential
contribution is the dominant force component.  For norm-conserving
pseudopotentials, the nonlocal contribution to the Hellmann--Feynman
force is non-zero and depends on the TD density via the same projectors
used in \texttt{CurrentNlpp.cpp}.  Active development on this lives in
the \texttt{develop-Ehrenfest} branch.  The current-operator analogue
$\ii[V_\text{NL}, \mathbf{r}]$ in \texttt{CurrentNlpp.cpp} is the
relevant reference; the force analogue is
$\frac{\partial}{\partial R_I} \langle \phi | V_\text{NL} | \phi \rangle$
contracted with $P(t)$.
\end{infobox}

% ============================================================================
\section{Forward look}
\label{sec:forward}
% ============================================================================

The Task~A/Task~B migration unblocks three lines of development:

\begin{enumerate}[leftmargin=*]
\item \textbf{Non-collinear spin RT-TDDFT.}  Requires fixing GAP-4
      (\texttt{vtot} must include the off-diagonal spin components of
      $V_{xc}$) in \texttt{tddft\_md()} l.~1039--1048.  The two existing
      \texttt{// noncoll need change} markers point to the exact lines.
      Also verify the bootstrap (Section~\ref{sec:bootstrap}) handles
      the $2\times 2$ spin block correctly.

\item \textbf{Ehrenfest dynamics with NL-PP forces.}  Open issue
      flagged in Section~\ref{sec:cve}.  The
      \texttt{develop-Ehrenfest} branch is the active worktree for this
      work; the relevant reference files in this branch are
      \texttt{Force/Force.cpp} (ground-state version) and
      \texttt{TDDFT/RMG\_TDDFT/CurrentNlpp.cpp} (current-operator
      analogue, annotated).

\item \textbf{Spin + Ehrenfest combined.}  Requires both of the above.
      The \texttt{rmg::tddft} class is already structured to support
      it: \texttt{tddft\_md()} is designed to be re-entrant from an
      outer MD loop, and the \texttt{first\_step} bootstrap rebuilds
      $H$ in the current geometry on every call.
\end{enumerate}

\paragraph{Pre-existing physics work that is still relevant.}
GAP-1 (sustained field) and GAP-2 (diamagnetic current) are the two
remaining items needed to extend RT-TDDFT from delta-kick optical-spectra
calculations to true pump--probe and laser-driven simulations.  Both
require changes in \texttt{rmg\_tddft.cpp} at the locations identified
in Section~\ref{sec:gaps}.

% ============================================================================
\section*{Document provenance}
% ============================================================================

\begin{itemize}[leftmargin=*]
  \item Frozen Rosetta Stone: \texttt{git show explore:TDDFT/RMG\_TDDFT/RmgTddft.cpp}
  \item Task~A commit (CurrentNlpp + VecPmatrix): \texttt{b01af2025}
  \item Task~B commit (rmg\_tddft.cpp annotations): \texttt{a998da7c2}
  \item Pre-sync anchor tag: \texttt{pre-TDDFT-annotation}
  \item Companion theory document: \texttt{docs/rmg\_tddft\_theory.tex}
        (in the frozen \texttt{explore} branch)
\end{itemize}

\end{document}
